\bookchapter{الفصل 21: تعمّق في شبكات الهاتف المحمول وأمن المعلومات}

\chapterintro{نستخدم هواتفنا يوميًا للتصفح والدفع والدردشة، ويبدو الأمر مريحًا للغاية. لكن هل فكرت يومًا في عدد المرات التي تكون فيها سجلات الدردشة وكلمات مرور الدفع وصورك “عارية” على الإنترنت؟ بصراحة، لم يفكر معظم الناس في هذه المشكلة أبدًا، لكن بمجرد وقوع الكارثة — مثل سرقة الحساب أو اختراق البطاقة المصرفية — فالأوان قد فات على الندم. في هذا الفصل سنتحدث جيدًا عن كيفية جعل حياتك الرقمية أقل “شفافية”.}

\booksection{أولًا: VPN — “عباءة الإخفاء” داخل الشبكة}

VPN، اختصار لـ Virtual Private Network (الشبكة الخاصة الافتراضية). يمكنك تخيلها كنفق تحت أرضي مخصص ومشفر بُني بين منزلك والموقع الذي تريد زيارته. تمشي داخل النفق ذهابًا وإيابًا، ولا يستطيع الناس في الخارج رؤية ما تحمله في يدك أو إلى أين أنت ذاهب.

\begin{itemize}
  \item \textbf{التصفح العادي}: هاتفك الشبكة العامة الموقع المستهدف. عبر أجهزة التوجيه والمحولات في المنتصف، يمكن لأي عقدة أن "تتلصص" على البيانات التي ترسلها. خاصة على شبكات Wi-Fi العامة في المقاهي والمطارات، بصراحة، الأمر أشبه بـ "العري".
  \item \textbf{استخدام VPN}: هاتفك نفق VPN المشفر خادم VPN الموقع المستهدف. تُغلَّف البيانات وتُشفَّر، ولا يرى "المتلصصون" سوى فوضى من الرموز. بالإضافة إلى ذلك، يرى الموقع المستهدف عنوان IP الخاص بخادم VPN وليس عنوان IP الحقيقي الخاص بك.
\end{itemize}

إذن ماذا يفعل VPN بالضبط؟ ببساطة: أولًا،\textbf{حماية خصوصيتك على شبكة Wi-Fi العامة}؛ ثانيًا،\textbf{تجاوز القيود الجغرافية}(مثل مشاهدة محتوى متاح في دول معينة فقط)؛ ثالثًا،\textbf{منع مزود خدمة الإنترنت (ISP) من تسجيل عادات التصفح الخاصة بك}. لكن يجب أن أحذرك:\textbf{لا تستخدم VPN مجاني!}فالخدمة المجانية غالبًا ما تعني أن بياناتك هي “سلعتهم”. لدي صديق استخدم VPN مجانيًا، وبدأ هاتفه فجأة يعرض إعلانات منبثقة، واكتشف لاحقًا أن ذلك الـ VPN باع بيانات التصفح الخاصة به لمعلنين. لذا عند اختيار VPN، إما أن تختار علامة تجارية مدفوعة كبيرة (مثل NordVPN أو ExpressVPN)، أو لا تستخدمه أبدًا.

بالإضافة إلى ذلك، VPN ليس حلًا سحريًا. فهو يحمي فقط البيانات “أثناء النقل”، وإذا كان الموقع المستهدف نفسه غير آمن أو كان هاتفك مصابًا ببرامج خبيثة، فلن يتمكن VPN من مساعدتك.

\booksection{ثانيًا: “الفخ الحلو” لشبكات Wi-Fi العامة — أخطر مما تتخيل}

قد تظن أن الأمر مجرد الاتصال بشبكة Wi-Fi مجانية، فما حجم الخطر؟ سأحكي لك قصة حقيقية: أجرى باحثون أمنيون تجربة، نصبوا فيها نقطة وصول Wi-Fi باسم “Starbucks\_Free” في المطار، وفي أقل من ساعة اتصل بها مئات الأشخاص. في الواقع، كانت هذه النقطة من إعداد الباحثين أنفسهم، وكل\textbf{حركة المرور غير المشفرة}التي تمر عبرها — مثل محتوى صفحات HTTP التي تزورها واتصالات التطبيقات التي لا تستخدم TLS — يمكنهم بالفعل اعتراضها وعرضها.

لكن انتبه لشرط أساسي: معظم تطبيقات WeChat والبنوك وAlPay التي تستخدمها يوميًا تستخدم بالفعل تشفير HTTPS. بصراحة، حتى لو اعترض المهاجم حزم البيانات، فلن يرى سوى فوضى من الرموز،\textbf{ولن يتمكن من قراءة محتوى الدردشة أو كلمة مرور الدفع مباشرة}. لكنهم سيظلون قادرين على معرفة المواقع التي زرتها (أسماء النطاقات)، وإذا كانت بعض التطبيقات أو الصفحات القديمة لا تزال تستخدم نقل HTTP بنص واضح، فسيكون المحتوى مكشوفًا بالكامل. يُسمى هذا النوع من الهجوم “هجوم الوسيط” (Man-in-the-Middle)، ومبدأه أن المهاجم “يضع قدمه” بينك وبين جهاز توجيه Wi-Fi الحقيقي. ما هو أكثر رعبًا أن هواتف الكثير من الناس مضبوطة على “الاتصال التلقائي بشبكات Wi-Fi المعروفة”، فإذا قام المهاجم بضبط اسم نقطة اتصاله على اسم شبكة Wi-Fi اتصلت بها سابقًا (مثل “ChinaNet” أو “CMCC”)، سيتصل هاتفك تلقائيًا.

إذن كيف تحمي نفسك؟ إليك بعض النصائح العملية:

\begin{itemize}
  \item \textbf{ابحث عن HTTPS}: بداية شريط العنوان هي https://، ويمثل "s" الإضافي كلمة Secure (آمن). عند زيارة أي موقع يتطلب إدخال كلمة مرور، إذا رأيت http:// (بدون s)، توقف فورًا ولا تدخل! الآن تعرض المتصفحات الرئيسية تحذير "غير آمن" قبل مواقع HTTPS غير المؤمنة، وعند رؤية هذه العلامة يجب الحذر.
  \item \textbf{أوقف "الاتصال التلقائي"}: سواء كان iPhone أو Android، أوقف خيار "الانضمام التلقائي" أو "الاتصال التلقائي" في إعدادات Wi-Fi. من الأفضل النقر يدويًا في كل مرة بدلًا من ترك الهاتف يقرر بنفسه.
  \item \textbf{استخدم بيانات الهاتف للعمليات الحساسة}: التحويلات المصرفية، تسجيل الدخول للبنوك، إدخال رقم الهوية... استخدم بيانات 4G/5G الخاصة بك بأمانة لهذه العمليات، ولا تتجنب تلك البيانات القليلة.
  \item \textbf{انسَ شبكات Wi-Fi غير المستخدمة}: عند العودة من رحلة عمل أو سفر، "انسَ" شبكات Wi-Fi الخاصة بالفنادق والمطارات التي اتصلت بها مؤقتًا. بهذه الطريقة لن يتصل هاتفك تلقائيًا بنقطة وصول مزيفة تحمل الاسم نفسه في المرة القادمة.
\end{itemize}

\booksection{ثالثًا: كيف يحميك HTTPS؟ — مصافحة SSL/TLS}

ذكرنا للتو HTTPS، فكيف يعمل بالضبط؟ في كل مرة تزور فيها موقع HTTPS، تجري بين المتصفح والخادم “مصافحة SSL/TLS”، والعملية برمتها تبدو هكذا تقريبًا:

\begin{enumerate}
  \item يقول متصفحك: "مرحبًا يا خادم، أريد التواصل معك بأمان، وهذه قائمة قدراتي التشفيرية (خوارزميات التشفير المدعومة وغيرها)."
  \item يرد الخادم: "حسنًا، لقد اخترت مجموعة التشفير هذه من قائمتك، وهذه شهادتي الرقمية التي تحتوي على مفتاحي العام، موقعة من جهة CA معتمدة."
  \item يتحقق المتصفح من صحة الشهادة: هل صدرت من CA موثوق؟ هل انتهت صلاحيتها؟ هل يطابق اسم النطاق؟
  \item بعد اجتياز التحقق، يتفاوض الطرفان باستخدام المفتاح العام لتشفير "مفتاح جلسة"، وبعدها يتم تشفير جميع الاتصالات بهذا المفتاح.
\end{enumerate}

هذا أشبه بذهابك إلى البنك: تقدم بطاقة هويتك (الشهادة) أولاً، يتحقق الموظف من هويتك، ثم تتفقان على شفرة سرية، وبعدها لا يفهم محادثتكما سواكما. حتى لو استمع شخص ما سرًا، فلن يسمع سوى رموز مشفرة غير مفهومة.

لكن انتبه: HTTPS يضمن فقط أن “عملية النقل” مشفرة، ولا يضمن أن الموقع المستهدف جيد. يمكن للمواقع الاحتيالية أيضًا استخدام HTTPS — أيقونة القفل الصغيرة تلك تثبت فقط أنك تتصل بـ “موقع ما”، ولا تعني أن ذلك الموقع هو الذي تريد الذهاب إليه. لذا عليك التحقق بنفسك من صحة عنوان الموقع.

\booksection{رابعًا: أمان DNS — هل “دليل هاتف الشبكة” الخاص بك آمن؟}

DNS (نظام أسماء النطاقات) يشبه دليل هاتف الإنترنت: عندما تكتب “baidu.com”، يساعدك DNS في العثور على عنوان IP المقابل (مثل 110.242.68.66)، ثم يتمكن المتصفح من الاتصال. المشكلة أن استعلامات DNS التقليدية\textbf{بنص واضح}و\textbf{غير مشفرة}! بمعنى أن مزود خدمة الإنترنت ومزود Wi-Fi وأي عقدة وسيطة يمكنها رؤية المواقع التي زرتها.

ما هو أسوأ من ذلك هو “اختطاف DNS”: يتلاعب المهاجم في مرحلة DNS لتوجيه “baidu.com” إلى موقع تصيد مزيف. أنت تكتب العنوان الصحيح، لكنك تُنقل إلى صفحة دخول مطابقة تمامًا ولكنها مزيفة في الواقع، وبمجرد إدخال حسابك وكلمة مرورك تنتهي القصة.

لحسن الحظ، أصبحت هناك الآن حلول DNS مشفرة:

\begin{itemize}
  \item \textbf{DNS over HTTPS (DoH)}: يخفي استعلامات DNS داخل حركة مرور HTTPS، ممزوجة مع حركة تصفح الويب العادية، فلا يمكن تمييز أيها استعلام DNS.
  \item \textbf{DNS over TLS (DoT)}: يستخدم بروتوكول تشفير TLS لحماية استعلامات DNS خصيصًا، بمنفذ مستقل، يمكن لمسؤول الشبكة معرفة أنك تستخدم DoT لكنه لا يرى المحتوى.
\end{itemize}

شخصيًا أوصي باستخدام 1.1.1.1 من Cloudflare أو 8.8.8.8 من Google كخوادم DNS، فكلاهما يدعم DNS المشفر. في إعدادات “DNS الخاص” على Android، ما عليك سوى إدخال “one.one.one.one” أو “dns.google”، وهو سهل للغاية. أما iPhone فأمره أصعب قليلًا، إذ يتطلب تثبيت ملف وصف أو الإعداد عبر تطبيق، لكنه يستحق العناء.

\booksection{خامسًا: هل “التصفح المتخفي” يمنحك إخفاءً حقيقيًا؟}

الخلاصة أولًا:\textbf{الوضع المتخفي هو ستارة نافذتك، وليس عباءة إخفاء.}

ما يفعله: لا يترك أثرًا محليًا — بعد إغلاق النافذة المتخفية، لا يحفظ المتصفح سجل التصفح أو ملفات تعريف الارتباط أو سجلات تعبئة النماذج أو سجلات البحث. استخدامه مناسب عند استعارة جهاز شخص آخر للبحث عن شيء ما، أو عند شراء هدية دون أن يكتشفها أفراد العائلة من سجل التصفح.

ما لا يستطيع فعله:

\begin{itemize}
  \item الموقع الذي تزوره\textbf{لا يزال يرى أنك جئت}، فعنوان IP الخاص بك ومعلومات جهازك وسلوك التصفح الخاص بك واضحة تمامًا لدى الموقع.
  \item مزود خدمة الإنترنت الخاص بك\textbf{لا يزال يرى المواقع التي زرتها}، لكنه لا يرى المحتوى المحدد (إذا استخدمت HTTPS).
  \item مسؤولو شبكة الشركة/المدرسة\textbf{قادرون أيضًا على المراقبة}لسلوكك على الإنترنت.
  \item الملفات التي تنزلها والإشارات المرجعية التي تضيفها\textbf{ستبقى محفوظة}، ولن تختفي عند إغلاق النافذة.
\end{itemize}

لذا تذكر: الوضع المتخفي هو “ممحاة محلية”، وليس “عباءة إخفاء شبكية”. إذا كنت بحاجة فعلية إلى إخفاء آثار تصفحك، فعليك استخدام تركيبة مثل VPN + متصفح Tor، لكن للاستخدام اليومي، معرفة هذه الأمور كافية.

\booksection{سادسًا: خصوصية المتصفح — لا تدع المتصفح يبيعك}

قد لا تدرك أن المتصفح بحد ذاته جامع ضخم للمعلومات. ملفات تعريف الارتباط، وسجل التصفح، وذاكرة التخزين المؤقت، والملء التلقائي، وأذونات الموقع... كل واحد منها قد يكشف خصوصيتك. بعض النصائح:

\begin{itemize}
  \item \textbf{تنظيف بيانات التصفح بانتظام}: كل فترة، اذهب إلى الإعدادات ونظف ملفات تعريف الارتباط وذاكرة التخزين المؤقت وسجل التصفح. ليس لمنع "التتبع" — فذلك لم يعد مجديًا بتنظيف ملفات تعريف الارتباط — بل لمنع الهاتف من أن يصبح أبطأ. يمكن أن تحتل ذاكرة التخزين المؤقت للمتصفح عدة غيغابايت.
  \item \textbf{تخلَّ عن وهم "عدم التتبع"}: خيار "إرسال طلب 'عدم التتبع'" في المتصفح هو في جوهره "طلب" من المواقع ألا تتبعك. لكن ما إذا كانت المواقع تستجيب هو شأن آخر تمامًا، ومعظم المواقع تتجاهله. لا تعتمد عليه.
  \item \textbf{استخدم متصفحات حماية الخصوصية}: Firefox Focus ومتصفح DuckDuckGo وBrave تحظر بشكل افتراضي أدوات التتبع والإعلانات، وهي أنظف بكثير من الإعدادات الافتراضية لـ Chrome وSafari.
  \item \textbf{محرك البحث مهم أيضًا}: استخدم DuckDuckGo بدلًا من Baidu/Google للبحث، فهو لا يسجل سجل بحثك ولا يبني ملفًا شخصيًا للمستخدم. بصراحة، نتائج البحث ليست دقيقة مثل Google، لكنها كافية للبحث اليومي.
  \item \textbf{تحقق من أذونات التطبيقات}: اذهب إلى إعدادات الهاتف لترى الأذونات الممنوحة لمتصفحك. هل يحتاج فعلًا إلى الوصول إلى موقعك وميكروفونك وجهات اتصالك؟ على الأرجح لا.
\end{itemize}

\booksection{سابعًا: قائمة أمان عملية — ابدأ بحماية نفسك من الآن}

بعد كل هذا الكلام، إليك قائمة أمان بسيطة ومباشرة، اقضِ نصف ساعة في إنجازها وستحدث نقلة نوعية في أمان هاتفك:

\begin{enumerate}
  \item \textbf{أوقف الاتصال التلقائي بـ Wi-Fi}: في iPhone من الإعدادات - شبكة لاسلكية، وفي Android من إعدادات Wi-Fi.
  \item \textbf{غيّر إلى DNS مشفر}: على Android اضبط DNS الخاص على "dns.google"، وعلى iPhone استخدم تطبيقًا أو ملف وصف لضبط DoH من 1.1.1.1.
  \item \textbf{فعّل المصادقة الثنائية (2FA)}: WeChat وAlPay وApple ID وحساب Google، فعّل 2FA لكل ما يمكن. رمز التحقق عبر الرسائل النصية هو أدنى مستوى من 2FA، واستخدام TOTP (مثل Google Authenticator أو Microsoft Authenticator) أفضل.
  \item \textbf{راجع أذونات التطبيقات}: مرر على أذونات كل تطبيق، خاصة الموقع وجهات الاتصال والرسائل والميكروفون والكاميرا. ستكتشف أن بعض تطبيقات المصباح تطلب جهات الاتصال — احذفها فورًا.
  \item \textbf{نظّف بيانات المتصفح بانتظام}: اضبط تذكيرًا في التقويم، مرة شهريًا، لتنظيف ذاكرة التخزين المؤقت وملفات تعريف الارتباط.
  \item \textbf{لا تستخدم كابلات البيانات في محطات الشحن العامة}: منفذ USB يشحن وينقل البيانات في الوقت ذاته، ويمكن لمحطات الشحن الخبيثة (المسماة "Juice Jacking") سرقة البيانات أثناء شحنك. احمل رأس الشحن والكابل الخاصين بك عند الخروج.
\end{enumerate}

\begin{tcbpractice}{افعل ذلك الآن: أجرِ "فحص أمان" لهاتفك}
\begin{enumerate}
  \item \textbf{ترقية DNS}: Android: الإعدادات > الاتصال والمشاركة > DNS الخاص > أدخل "one.one.one.one" أو "dns.google". iPhone: ابحث في Safari عن "1.1.1.1 DNS profile" لتثبيت ملف الوصف، أو نزّل تطبيق 1.1.1.1 من Cloudflare.
  \item \textbf{تدقيق Wi-Fi}: افتح قائمة Wi-Fi، و"انسَ" النقاط غير المعروفة وغير المستخدمة وذات الأسماء الغريبة.
  \item \textbf{تنظيف الأذونات}: ادخل إلى الإعدادات > الخصوصية/إدارة الأذونات، وراجع أذونات كل تطبيق. حدد التطبيقات ذات طلبات الأذونات غير المعقولة (مثل آلة حاسبة تطلب جهات الاتصال)، وفكر في حذفها أو استخدام بدائل.
  \item \textbf{تحقق من أمان الحساب}: افتح صفحة أمان Apple ID أو حساب Google، لترى ما إذا كان هناك تسجيل دخول من أجهزة غير معروفة، وتأكد من تفعيل 2FA، وتحقق من موثوقية جهات اتصال الاستعادة.
\end{enumerate}

\end{tcbpractice}
